\section{EMC Cost Categorization}\label{emc-cost-categorization-1}

EMC has two assigned cost categories, referred to as \emph{standard} and \emph{product}. These two categories are further divided into \emph{software} and \emph{hardware}. For both Centera and CLARiiON, \emph{standard cost} results in a \textbf{10\% uplift} to (hardware) material cost; there is no uplift assigned to the C* software.

\emph{Standard margin} is reported against \emph{standard cost} at the division level. This is the GM (gross margin) figure reported at the General Manager level and the one we are concerned about in our model. \emph{Product cost} is an accounting number presented to field Sales, principally as a means of setting a discount floor. It differs -- somewhat significantly -- by product. However, we do not have to consider \emph{product cost} in our calculations.

Figure~\ref{fig:cost} below illustrates this cost categorization.


\begin{figure}[htbp]
\centering
\resizebox{.85\linewidth}{!}{% Editable cost-category hierarchy, preserving the historical relationships.
\begin{tikzpicture}[x=1cm,y=1cm,capacity text]
\node[cost box,minimum width=2cm] (material) at (0,0) {Material\\cost};
\node[cost box] (software) at (4,3) {Software\\\textcolor{capacityblue}{Standard cost}};
\node[cost box] (hardware) at (4,1.35) {Hardware\\\textcolor{capacityblue}{Standard cost}};
\node[cost box] (product) at (4,-1.35) {Hardware\\\textcolor{capacitygreen}{Product cost}};
\node[cost box] (list) at (4,-3) {List price};
\foreach \target in {software,hardware,product,list}
  \draw[capacity axis] (material.east) -- ++(.35,0) |- (\target.west);
\draw[capacityred,dashed] (1.4,0) -- (8.4,0);
\draw[decorate,decoration={brace,amplitude=7pt},capacity axis]
  (5.9,3.6) -- (5.9,.75) node[midway,right=12pt,align=center] {Standard\\margin};
\draw[decorate,decoration={brace,amplitude=7pt},capacity axis]
  (5.9,-.75) -- (5.9,-3.6) node[midway,right=12pt,align=center] {Field\\margin};
\end{tikzpicture}
}

\caption{EMC Cost Categorization}\label{fig:cost}
\end{figure}
